% from ega4-11-fr.tex
\oldpage[IV-3]{126}

\item Il existe un schéma affine noethérien $S_0=\operatorname{Spec}(A_0)$,
un morphisme de type fini $f_0:X_0\to S_0$, un
$\mathcal O_{X_0}$-Module cohérent $\mathcal F_0$, un morphisme $S\to S_0$
tels que le $S$-préschéma $X_0\otimes_{S_0}S$ soit $S$-isomorphe à $X$ et
que, si on identifie $X$ à $X_0\otimes_{S_0}S$, $\mathcal F$ soit isomorphe
à $\mathcal F_0\otimes_{\mathcal O_{S_0}}\mathcal O_S$, et $\mathcal F_0$
soit $f_0$-plat au point $x_0$ projection de $x$ dans $X_0$ (resp.
$f_0$-plat).
\item Les conditions de b) sont vérifiées et en outre $A_0$ est une
sous-$\mathbf Z$-algèbre de type fini de $A$, le morphisme $S\to S_0$
correspondant à l'injection canonique $A_0\to A$.
\end{enumerate}
\end{corollary}

Il est clair que c) implique b) et b) implique a) en vertu de (2.1.4).
D'autre part, on peut considérer $A$ comme limite inductive de ses
sous-$\mathbf Z$-algèbres de type fini, et on sait par (8.9.1) qu'il y a une
telle sous-algèbre $A_0$ et un morphisme de type fini $f_0:X_0\to S_0$ tels
que $X$ soit $S$-isomorphe à $X_0\times_{S_0}S$ et $\mathcal F$ isomorphe à
$\mathcal F_0\otimes_{\mathcal O_{S_0}}\mathcal O_S$~; on peut donc écrire
$X=\varprojlim X_\lambda$, où $X_\lambda=X_0\otimes_{A_0}A_\lambda$,
$A_\lambda$ parcourant l'ensemble des sous-$\mathbf Z$-algèbres de type
fini de $A$ contenant $A_0$, et $\mathcal F=v_\lambda^*(\mathcal F_\lambda)$,
où $v_\lambda:X\to X_\lambda$ est la projection canonique et
$\mathcal F_\lambda=\mathcal F_0\otimes_{A_0}A_\lambda$. Il résulte de
(11.2.6) qu'il existe $\lambda$ tel que $\mathcal F_\lambda$ soit
$f_\lambda$-plat au point $x_\lambda=v_\lambda(x)$ (resp.
$f_\lambda$-plat)~; alors le sous-anneau $A_\lambda$ de $A$ vérifie les
conditions de c).

\begin{proposition}[11.2.8]
\label{IV.11.2.8-fr}
Soient $g:X\to S$, $h:Y\to S$ deux morphismes de présentation finie,
$f:X\to Y$ un $S$-morphisme, $\mathcal F$ un
$\mathcal O_X$-Module quasi-cohérent de présentation finie~; pour tout
$s\in S$, soient $X_s=g^{-1}(s)=X\times_S\operatorname{Spec}(k(s))$,
$Y_s=h^{-1}(s)=Y\times_S\operatorname{Spec}(k(s))$, $f_s$ le morphisme
$f\times1:X_s\to Y_s$, $\mathcal F_s=\mathcal F\otimes_{\mathcal O_S}k(s)$.
Alors l'ensemble $E$ des $s\in S$ tels que $\mathcal F_s$ soit $f_s$-plat
est localement constructible.
\end{proposition}

La propriété dont nous voulons démontrer qu'elle est constructible vérifie la
condition (9.2.1, (i)), en vertu de (2.2.13) et (2.5.1). Compte tenu de
(9.2.3), on peut donc se limiter au cas où $S$ est affine, noethérien et
intègre de point générique $\eta$ et à prouver que $E$ ou $S-E$ est un
voisinage de $\eta$ dans $S$. Si $\eta\in S-E$, cela résulte aussitôt du
lemme (9.4.7.1). Si au contraire $\eta\in E$, il résulte tout d'abord de
(11.2.6), appliqué suivant la méthode de (8.1.2, a)), qu'il y a un voisinage
ouvert $U$ de $\eta$ dans $S$ tel que
$\mathcal F|g^{-1}(U)$ soit $h^{-1}(U)$-plat~; \emph{a fortiori} (2.1.4)
$\mathcal F_s$ est $f_s$-plat pour tout $s\in U$, ce qui achève la
démonstration.

Le théorème suivant généralise (11.2.6.1, (ii))~:

\begin{theorem}[11.2.9]
\EGAReaderTarget{ega4.c04.thm.11-2-9}{IV.11.2.9-fr}
(Raynaud). --- Soient $A_0$ un anneau, $(A_\lambda)_{\lambda\in L}$ un système
inductif filtrant de $A_0$-algèbres, $B_0$ une $A_0$-algèbre de présentation
finie, $\mathfrak J_0$ un idéal de type fini de $B_0$, $M_0$ un $B_0$-module
de présentation finie. On pose
$B_\lambda=B_0\otimes_{A_0}A_\lambda$,
$\mathfrak J_\lambda=\mathfrak J_0B_\lambda$,
$M_\lambda=M_0\otimes_{A_0}A_\lambda=M_0\otimes_{B_0}B_\lambda$,
$A=\varinjlim A_\lambda$,
$B=\varinjlim B_\lambda=B_0\otimes_{A_0}A$,
$\mathfrak J=\mathfrak J_0B$,
$M=\varinjlim M_\lambda=M_0\otimes_{A_0}A=M_0\otimes_{B_0}B$. Pour que
$\operatorname{gr}_{\mathfrak J}^*(M)$ soit un $A$-module plat, il faut et il
suffit qu'il existe $\lambda$ tel que
$\operatorname{gr}_{\mathfrak J_\lambda}^*(M_\lambda)$ soit un
$A_\lambda$-module plat~; les homomorphismes canoniques
\begin{equation}
\left\{
\begin{aligned}
\operatorname{gr}_{\mathfrak J_\lambda}^*(M_\lambda)
  \otimes_{A_\lambda}A&\longrightarrow
  \operatorname{gr}_{\mathfrak J}^*(M),\\
\operatorname{gr}_{\mathfrak J_\lambda}^*(M_\lambda)
  \otimes_{A_\lambda}A_\mu&\longrightarrow
  \operatorname{gr}_{\mathfrak J_\mu}^*(M_\mu)
\end{aligned}
\right.
\qquad\text{pour }\lambda\leq\mu
\label{IV.11.2.9.1-fr}
\tag{11.2.9.1}
\end{equation}
sont alors bijectifs et
$\operatorname{gr}_{\mathfrak J_\mu}^*(M_\mu)$ est un $A_\lambda$-module plat
pour $\lambda\leq\mu$.
\end{theorem}

\begin{proof}
Montrons d'abord que les conditions sont suffisantes, ce qui revient à prouver
la bijectivité de (11.2.9.1). Cela résulte du lemme suivant~:

\begin{lemma}[11.2.9.2]
\label{IV.11.2.9.2-fr}
Soient $A$ un anneau, $B$ une $A$-algèbre, $M$ un $B$-module, $\mathfrak J$ un
idéal de $B$. Soient $A'$ une $A$-algèbre~; posons
$B'=B\otimes_AA'$, $M'=M\otimes_AA'=M\otimes_BB'$,
$\mathfrak J'=\mathfrak JB'$. Si
$\operatorname{gr}_{\mathfrak J}^*(M)$ est un $A$-module plat,
l'homomorphisme canonique
\[
(\operatorname{gr}_{\mathfrak J}^*(M))\otimes_AA'
\longrightarrow \operatorname{gr}_{\mathfrak J'}^*(M')
\]
est bijectif.
\end{lemma}

